\documentclass[11pt,a4paper]{article}

%================================================
% Packages
%================================================

\usepackage[utf8]{inputenc}
\usepackage[T1]{fontenc}

\usepackage{graphicx}
\usepackage{amsmath}
\usepackage{amssymb}
\usepackage{booktabs}
\usepackage{longtable}
\usepackage{array}
\usepackage{enumitem}
\usepackage{geometry}
\usepackage{titlesec}
\usepackage{xcolor}
\usepackage{microtype}
\usepackage[section]{placeins}
\usepackage{hyperref}

%================================================
% Page layout
%================================================

\geometry{
    top=2.5cm,
    bottom=2.5cm,
    left=2.5cm,
    right=2.5cm
}

\raggedbottom

%================================================
% Hyperlinks
%================================================

\hypersetup{
    colorlinks=true,
    linkcolor=blue,
    citecolor=blue,
    urlcolor=blue
}

%================================================
% Heading depth
%================================================

% Number headings down to subsubsection level
\setcounter{secnumdepth}{3}

% Include headings down to subsubsection level in the table of contents
\setcounter{tocdepth}{3}

% Display paragraph headings on their own line
\titleformat{\paragraph}[block]
    {\normalfont\normalsize\bfseries}
    {\theparagraph}
    {1em}
    {}

\titlespacing*{\paragraph}
    {0pt}
    {3.25ex plus 1ex minus 0.2ex}
    {1.2ex plus 0.2ex}

%================================================
% Compact lists
%================================================

\setlist[itemize]{
    leftmargin=2em,
    itemsep=0.2em,
    topsep=0.5em,
    parsep=0pt
}

%================================================
% Title
%================================================

\title{
    \textbf{Artificial Intelligence for Earth Observation}\\
    \vspace{0.3cm}
    \large An Overview of the Current Research Landscape and Open Challenges
}

\author{Stefano Infusini}
\date{August 2026}

%================================================
% Document
%================================================

\begin{document}

\maketitle

\tableofcontents

\newpage

%================================================
\section{Introduction}

Earth Observation (EO) provides systematic information about the Earth's surface, oceans, and atmosphere through a variety of remote-sensing technologies, including optical imaging (RGB, multispectral, and hyperspectral), radar, and LiDAR sensors. The increasing availability of satellite missions and the continuous improvement of spatial, spectral, and temporal resolution have resulted in a rapidly growing volume of EO data.

Artificial Intelligence (AI) has consequently become an increasingly important component of Earth Observation. Machine-learning and deep-learning approaches are now applied to a wide range of tasks, including scene classification, semantic segmentation and land-cover mapping, object detection, change detection, geophysical parameter estimation, and environmental forecasting. More recently, research has expanded toward self-supervised learning, multimodal learning, Transformer-based architectures, and large-scale foundation models for EO.

This report provides a technical overview of the main concepts and current research directions at the intersection of Artificial Intelligence and Earth Observation. It first introduces the EO data modalities and remote-sensing characteristics needed to understand the field, and then focuses on the principal AI4EO tasks, the evolution toward modern representation learning, current research trends, and the main open challenges. The report is intended to provide a common basis for discussing the scope and feasibility of possible thesis directions, rather than to preselect a specific thesis topic.


%================================================
\section{Earth Observation Fundamentals}

\subsection{What is Earth Observation?}

Earth Observation (EO) is the process of acquiring information about
the Earth's surface, oceans, and atmosphere using remote-sensing
instruments \cite{esa_earth_observation}. These instruments can be
deployed on satellites, aircraft, unmanned aerial vehicles, or
ground-based platforms. Satellite observations are particularly
important because they provide systematic and repeated measurements
over large and potentially inaccessible geographical areas.

Remote-sensing instruments can be divided into passive and active
sensors. Passive sensors measure naturally available electromagnetic
radiation, such as reflected sunlight or thermal radiation emitted by
the Earth. Active sensors, such as Synthetic Aperture Radar and LiDAR,
transmit their own electromagnetic signal and measure the returned
response.

The acquired measurements are processed into products such as optical
images, radar images, elevation models, atmospheric measurements, and
geophysical variables. EO data are used for applications including
environmental monitoring, agriculture, climate research, disaster
management, urban planning, ocean observation, and land-cover mapping.

\subsection{EO Data Modalities}

Optical remote sensing is a broad family that includes conventional RGB,
multispectral, and hyperspectral imaging. These cases are discussed
separately below because they differ substantially in the number, width,
and placement of the spectral bands that are acquired. Radar and LiDAR,
by contrast, are active sensing modalities based on microwave and laser
measurements, respectively.

%------------------------------------------------
\subsubsection{Optical Imagery}

In Earth Observation, optical imagery consists of images of the Earth's
surface acquired by sensors that measure reflected solar radiation.
Optical satellite sensors operate according to a principle similar to
that of conventional cameras; however, they often acquire information
over a wider range of wavelengths than the three visible colour channels
used in standard RGB imagery.

Optical instruments are generally nadir-viewing or near-nadir-viewing
sensors, with spatial resolutions ranging approximately from
\(1\,\mathrm{m}\) to \(300\,\mathrm{m}\) and swath widths in the order
of tens to hundreds of kilometres \cite{esa_optical_imagery}.

The acquisition process can be represented schematically as follows:

\[
\text{Sun}
\longrightarrow
\text{Earth's surface}
\longrightarrow
\text{Satellite sensor}.
\]

Solar radiation reaching the Earth's surface is partially absorbed and
partially reflected. Different materials reflect different amounts of
electromagnetic radiation at different wavelengths. Consequently,
vegetation, water, soil, snow, rock, and artificial surfaces exhibit
distinct spectral responses, commonly referred to as
\textit{spectral signatures}.

\begin{figure}[!htbp]
    \centering

    \IfFileExists{images/reflectance_plot.jpg}{
        \includegraphics[
            width=\linewidth,
            height=0.42\textheight,
            keepaspectratio
        ]{images/reflectance_plot.jpg}
    }{
        \fbox{
            \parbox[c][5cm][c]{0.9\linewidth}{
                \centering
                Image not found.\\[0.3cm]
                Add the file as:\\
                \texttt{images/reflectance\_plot.jpg}
            }
        }
    }

    \caption{
        Representative spectral reflectance curves of snow, rock,
        vegetation, and water across the visible and near-infrared
        (VNIR), short-wave infrared (SWIR), and thermal infrared
        (TIR) regions. The figure also illustrates atmospheric
        transmission and the spectral bands and spatial resolutions
        of the ASTER and Landsat ETM+ sensors. Source: ESA
        \cite{esa_reflectance_curves}.
    }

    \label{fig:reflectance-curves-earth-features}
\end{figure}

As shown in Figure~\ref{fig:reflectance-curves-earth-features}, different
surface materials exhibit distinct wavelength-dependent behaviours.
Healthy vegetation absorbs a large proportion of visible blue and red
radiation because of photosynthetic pigments, while reflecting more
green radiation. Its reflectance then increases sharply between the red
and near-infrared regions, producing the feature known as the
\textit{red edge}. Vegetation exhibits high near-infrared reflectance
because of scattering within the internal structure of plant leaves.

Water generally has low reflectance, particularly in the near-infrared
and short-wave infrared regions, where radiation is strongly absorbed.
Snow is highly reflective in the visible region, but its reflectance
decreases considerably towards the short-wave infrared. Rock and bare
soil generally exhibit smoother spectral responses, whose precise shapes
depend on properties such as mineral composition, moisture, and surface
roughness.

The atmospheric-transmission curve in
Figure~\ref{fig:reflectance-curves-earth-features} also shows that the
atmosphere does not transmit electromagnetic radiation uniformly across
all wavelengths. Optical sensors are therefore designed to acquire
measurements within relatively transparent regions of the spectrum,
known as \textit{atmospheric windows}.

The satellite sensor records the spectral radiance reaching the
instrument within a set of predefined spectral bands. After radiometric
calibration and, when required, atmospheric correction, these
measurements can be converted into reflectance values. Each image pixel
therefore contains a set of spectral values describing the corresponding
area of the Earth's surface.

For example, a Sentinel-2 pixel \(p\) can be represented as a
13-dimensional spectral vector:

\[
\mathbf{x}_p =
\left[
B_1,
B_2,
B_3,
B_4,
B_5,
B_6,
B_7,
B_8,
B_{8A},
B_9,
B_{10},
B_{11},
B_{12}
\right].
\]

Although the numbering ends at \(B_{12}\), Sentinel-2 acquires 13
spectral bands because \(B_{8A}\) is an additional narrow
near-infrared band.

A true-colour Sentinel-2 image assigns the visible red, green, and blue
bands to the corresponding display channels:

\[
\mathrm{RGB}_{\mathrm{true}}
=
[B_4,B_3,B_2].
\]

A false-colour image assigns one or more non-visible bands to the RGB
display channels. A common false-colour combination is:

\[
\mathrm{RGB}_{\mathrm{false}}
=
[B_8,B_4,B_3].
\]

In this combination, the near-infrared band \(B_8\) is displayed through
the red channel. Healthy vegetation therefore generally appears bright
red because it strongly reflects near-infrared radiation.

%------------------------------------------------
\paragraph{Applications and Limitations}

Optical imagery is widely used to monitor vegetation and crop health,
forests, water bodies and floods, burned areas, urban development, snow
and ice, land-cover changes, and coastal environments. These applications
are possible because different materials exhibit distinct spectral
signatures and therefore reflect electromagnetic radiation differently
across wavelengths.

Most optical satellite sensors are passive sensors, meaning that they
depend on an external source of illumination, usually sunlight.
Consequently, optical observations can be affected by clouds, shadows,
haze, atmospheric conditions, and the absence of sunlight at night.

This represents one of the main differences between optical and radar
imagery. Active radar sensors, such as the Synthetic Aperture Radar
(SAR) instrument carried by Sentinel-1, transmit microwave signals
towards the Earth's surface and measure the returned signal. Radar
systems can therefore acquire observations both during the day and at
night and can generally operate through clouds.

%------------------------------------------------
%------------------------------------------------
\subsubsection{Multispectral Imagery}

Multispectral imagery is acquired by measuring electromagnetic radiation
within a finite number of discrete wavelength intervals, referred to as
\textit{spectral bands}. Unlike a conventional RGB image, which contains
only red, green, and blue channels, a multispectral image may also contain
bands in the red-edge, near-infrared, and short-wave infrared regions.

After the bands have been geometrically aligned and resampled to a common
spatial grid, a multispectral image can be represented as a three-dimensional
data cube:

\[
\mathbf{X}
\in
\mathbb{R}^{H \times W \times B},
\]

where \(H\) and \(W\) represent the spatial dimensions of the image and
\(B\) is the number of spectral bands. A pixel \(p\) is therefore
represented by a spectral vector:

\[
\mathbf{x}_p =
\left[
x_{p,1},
x_{p,2},
\ldots,
x_{p,B}
\right]
\in
\mathbb{R}^{B}.
\]

Each component \(x_{p,b}\) corresponds to the measurement acquired in
spectral band \(b\). A band does not measure a single wavelength.
Instead, it integrates the radiation reaching the sensor over a finite
wavelength interval. This process can be approximated as:

\[
x_{p,b}
=
\frac{
    \displaystyle
    \int L_p(\lambda)S_b(\lambda)\,d\lambda
}{
    \displaystyle
    \int S_b(\lambda)\,d\lambda
},
\]

where \(L_p(\lambda)\) is the spectral radiance arriving from pixel \(p\)
and \(S_b(\lambda)\) is the spectral response function of band \(b\).
The spectral response function describes the sensitivity of the sensor
to different wavelengths within the band.

Multispectral sensors therefore provide a discrete and relatively coarse
sampling of the continuous spectral signature of a surface. Their bands
are deliberately positioned in wavelength regions that are informative
for particular Earth-surface properties and that generally correspond to
atmospheric transmission windows.

For example, visible bands provide information about surface colour and
photosynthetic pigments, red-edge bands are sensitive to vegetation
condition, near-infrared bands are useful for vegetation structure and
land--water discrimination, and short-wave infrared bands provide
information about moisture, burned surfaces, snow, ice, and minerals.

The Sentinel-2 Multispectral Instrument is an important example of a
multispectral Earth Observation sensor. It acquires 13 spectral bands,
with four bands at a spatial resolution of \(10\,\mathrm{m}\), six at
\(20\,\mathrm{m}\), and three at \(60\,\mathrm{m}\)
\cite{esa_sentinel2_msi}. Consequently, the original Sentinel-2 bands do
not all represent ground areas of the same size.

To construct a single tensor containing all bands, the coarser bands are
commonly resampled to a shared spatial resolution. However, resampling a
\(20\,\mathrm{m}\) or \(60\,\mathrm{m}\) band to \(10\,\mathrm{m}\)
does not create new spatial information. It only interpolates the
original measurements onto a finer grid.

Multispectral bands can also be combined mathematically to emphasize
particular surface properties. For example, the Normalized Difference Vegetation Index (NDVI) is defined as
\cite{rouse1974monitoring}:

\[
\mathrm{NDVI}
=
\frac{\rho_{\mathrm{NIR}}-\rho_{\mathrm{Red}}}
     {\rho_{\mathrm{NIR}}+\rho_{\mathrm{Red}}}.
\]

where \(\rho_{\mathrm{NIR}}\) and \(\rho_{\mathrm{Red}}\) denote
reflectance in the near-infrared and red bands, respectively. Healthy
vegetation generally absorbs red radiation and strongly reflects
near-infrared radiation, producing relatively high NDVI values.

\paragraph{Applications and Limitations}

Multispectral imagery is widely used for land-cover classification,
vegetation monitoring, precision agriculture, water detection, flood
mapping, wildfire assessment, snow monitoring, urban analysis, and
change detection. Its main advantage over RGB imagery is the ability to
observe surface properties that are not directly visible to the human
eye.

Nevertheless, multispectral sensors measure only a limited number of
relatively broad bands. Narrow absorption features may therefore be
lost, and materials with similar responses in the available bands may
remain difficult to distinguish. Multispectral observations are also
affected by atmospheric interference, clouds, shadows, differences in
spatial resolution between bands, mixed pixels, and variations in
illumination and viewing geometry.

%------------------------------------------------
\subsubsection{Hyperspectral Imagery}

Hyperspectral imagery, also referred to as
\textit{imaging spectroscopy}, measures the spectrum of the radiation
reaching the sensor for every spatial element in an image
\cite{jpl_imaging_spectroscopy}. Like multispectral imagery, it can be
represented as a three-dimensional data cube:

\[
\mathbf{X}
\in
\mathbb{R}^{H \times W \times B}.
\]

However, in hyperspectral imagery, \(B\) is generally much larger.
Hyperspectral sensors typically acquire tens or hundreds of narrow,
closely spaced, and often contiguous spectral bands. A pixel is therefore
represented by a detailed spectral vector:

\[
\mathbf{x}_p =
\left[
\rho_p(\lambda_1),
\rho_p(\lambda_2),
\ldots,
\rho_p(\lambda_B)
\right],
\]

where \(\rho_p(\lambda_b)\) represents the reflectance measured around
wavelength \(\lambda_b\).

The spectral resolution of a band can be described through its bandwidth:

\[
\Delta\lambda_b
=
\lambda_{b,\max}
-
\lambda_{b,\min}.
\]

A smaller value of \(\Delta\lambda_b\) indicates a narrower band and
therefore a higher ability to distinguish fine spectral features.
Spectral resolving power may also be expressed as:

\[
R_b =
\frac{\lambda_b}{\Delta\lambda_b},
\]

where \(\lambda_b\) is the central wavelength of the band.

There is no universal numerical threshold that completely separates
multispectral from hyperspectral imagery. The main distinction is that
multispectral sensors observe selected regions using a relatively small
number of broader bands, whereas hyperspectral sensors densely sample a
spectral interval using many narrow and approximately contiguous bands.

This dense spectral sampling allows hyperspectral sensors to capture
narrow absorption features associated with the physical and chemical
composition of materials. Consequently, materials that appear similar
in RGB or multispectral imagery may become distinguishable through their
detailed hyperspectral signatures.

The NASA Airborne Visible/Infrared Imaging Spectrometer, for example,
samples the solar-reflected spectrum between approximately
\(400\,\mathrm{nm}\) and \(2500\,\mathrm{nm}\) using 224 spectral
channels \cite{jpl_imaging_spectroscopy}. Such measurements can support
the identification of minerals, vegetation constituents, water
properties, atmospheric components, and other surface materials.

Because the spatial resolution of remote-sensing imagery is finite, a
single pixel may contain multiple materials. In a simplified linear
spectral-mixture model, an observed pixel can be represented as:

\[
\mathbf{x}_p
=
\sum_{k=1}^{K}
a_{p,k}\mathbf{s}_k
+
\boldsymbol{\epsilon}_p,
\]

subject to:

\[
a_{p,k}\geq 0,
\qquad
\sum_{k=1}^{K}a_{p,k}=1.
\]

Here, \(\mathbf{s}_k\) is the spectral signature of material \(k\),
known as an \textit{endmember}, \(a_{p,k}\) is its estimated fractional
abundance within pixel \(p\), and \(\boldsymbol{\epsilon}_p\) represents
noise and modelling errors. The process of estimating endmembers and
their abundances is known as \textit{spectral unmixing}.

The large number of hyperspectral bands provides detailed information,
but also introduces substantial redundancy because neighbouring bands
are often strongly correlated. Hyperspectral datasets are therefore
frequently processed using band selection or dimensionality-reduction
methods such as Principal Component Analysis, Minimum Noise Fraction
transforms, autoencoders, or learned feature extractors.

For machine-learning applications, hyperspectral information can be
analysed spectrally, spatially, or jointly. One-dimensional convolutions
can learn patterns along the spectral dimension, two-dimensional
convolutions can extract spatial structures, and three-dimensional
convolutions can jointly process spectral and spatial information.

\paragraph{Applications and Limitations}

Hyperspectral imagery is used for mineral identification, precision
agriculture, vegetation-species discrimination, plant-stress analysis,
soil characterization, water-quality monitoring, environmental
contamination assessment, target detection, and spectral unmixing. Its
principal advantage is its ability to distinguish materials through
fine spectral characteristics that may not be captured by multispectral
sensors.

Its main limitations are its high dimensionality, large storage and
computational requirements, sensitivity to sensor noise and atmospheric
effects, and the limited availability of labelled training data.
Furthermore, some wavelength regions are strongly affected by atmospheric
water absorption and may need to be removed before analysis.

The combination of high dimensionality and limited training samples can
also reduce classification performance, since increasing the number of
features does not necessarily improve generalization. Dimensionality
reduction, regularization, self-supervised learning, and spectral-spatial
models are therefore particularly important in hyperspectral image
analysis \cite{bioucasdias2013hyperspectral}.

\subsubsection{Synthetic Aperture Radar}

Synthetic Aperture Radar (SAR) is an active microwave imaging technology
used to observe the Earth's surface. Unlike optical sensors, which depend
on reflected solar radiation, a SAR instrument transmits electromagnetic
pulses toward the ground and measures the portion of the signal scattered
back toward the sensor, known as \textit{backscatter}
\cite{moreira2013sar}.

The acquisition process can therefore be represented as

\[
\text{SAR sensor}
\longrightarrow
\text{Earth's surface}
\longrightarrow
\text{SAR sensor}.
\]

Because SAR systems provide their own source of illumination, they can
acquire images during both day and night. Moreover, the microwave
wavelengths used by SAR can generally propagate through clouds, making
radar imagery particularly useful when optical observations are
unavailable because of cloud cover or insufficient solar illumination
\cite{esa_sentinel1_instrument}.

SAR instruments are generally side-looking rather than nadir-viewing.
The image is consequently described using two principal directions:
\textit{range}, which is approximately perpendicular to the flight
direction, and \textit{azimuth}, which is parallel to the satellite
trajectory. The distance \(R\) between the sensor and a target can be
estimated from the round-trip travel time of the transmitted pulse:

\[
R = \frac{c\Delta t}{2},
\]

where \(c\) is the speed of light and \(\Delta t\) is the time elapsed
between transmission and reception of the radar signal.

The term \textit{synthetic aperture} refers to the method used to obtain
fine spatial resolution in the azimuth direction. As the satellite moves
along its orbit, the same target is observed from multiple successive
positions. The amplitude and phase of these observations are coherently
combined, making the system behave as if the measurements had been
acquired using an antenna much longer than the physical antenna carried
by the satellite. This virtual antenna is known as the synthetic
aperture \cite{moreira2013sar}.

A focused Single Look Complex (SLC) SAR pixel is represented by a complex
number:

\[
S(p) = I(p) + jQ(p) = A(p)e^{j\phi(p)},
\]

where \(A(p)\) is the amplitude of the returned signal and \(\phi(p)\) is
its phase. The corresponding intensity is

\[
I_{\mathrm{SAR}}(p) = |S(p)|^2.
\]

The amplitude or intensity describes the strength of the received echo,
whereas the phase contains information related to the distance travelled
by the electromagnetic wave. In detected products, such as Sentinel-1
Ground Range Detected (GRD) products, the complex phase information is
not preserved. After radiometric calibration, the measured intensity is
commonly expressed using a backscatter coefficient such as
\(\sigma^0\), often represented in decibels:

\[
\sigma^0_{\mathrm{dB}}
=
10\log_{10}\left(\sigma^0\right).
\]

The strength of the radar backscatter depends on several properties of
both the observed surface and the acquisition system:

\[
\sigma^0 =
f\left(
\lambda,
\theta,
P,
\varepsilon,
r,
g
\right),
\]

where \(\lambda\) is the radar wavelength, \(\theta\) is the incidence
angle, \(P\) is the polarization, \(\varepsilon\) represents the
dielectric properties of the target, \(r\) represents surface roughness,
and \(g\) represents the observation geometry.

Surface roughness must be considered relative to the radar wavelength.
A surface that is smooth compared with the wavelength produces mainly
specular reflection, directing most of the transmitted energy away from
the sensor. Calm water therefore generally produces low backscatter and
appears dark in a SAR image. Rough surfaces scatter energy in several
directions, including back toward the sensor, and consequently tend to
appear brighter.

Four principal scattering mechanisms are commonly considered:

\begin{itemize}
    \item \textit{surface scattering}, produced by interaction with a
    relatively rough surface such as bare soil;

    \item \textit{specular reflection}, produced by smooth surfaces such
    as calm water, which generally results in low backscatter;

    \item \textit{volume scattering}, produced by multiple interactions
    within a three-dimensional medium such as a vegetation canopy;

    \item \textit{double-bounce scattering}, produced by successive
    reflections between approximately perpendicular surfaces, such as
    the ground and a building wall or the ground and a tree trunk.
\end{itemize}

Moisture also strongly influences the radar response because water
changes the dielectric properties of soil and vegetation. Wet soil
therefore often produces stronger backscatter than dry soil, although
the final response also depends on surface roughness, incidence angle
and vegetation cover.

SAR signals can be transmitted and received using horizontal or vertical
polarization. The principal polarization channels are

\[
HH,\qquad HV,\qquad VH,\qquad VV,
\]

where the first letter represents the transmitted polarization and the
second represents the received polarization. The \(HH\) and \(VV\)
channels are referred to as co-polarized channels, whereas \(HV\) and
\(VH\) are cross-polarized channels. Different polarization combinations
are sensitive to different scattering mechanisms. For example,
cross-polarized backscatter is often useful for characterising complex
vegetation structures.

A Sentinel-1 observation acquired using the \(VV\) and \(VH\)
polarizations may therefore be represented as

\[
\mathbf{x}_p =
\left[
\sigma^0_{VV}(p),
\sigma^0_{VH}(p)
\right].
\]

SAR imagery also presents characteristics that complicate its
interpretation. Because radar is a coherent imaging system, SAR images
are affected by \textit{speckle}, a granular pattern caused by the
constructive and destructive interference of the echoes produced by
multiple scatterers within the same resolution cell. Speckle is an
inherent property of coherent radar imaging rather than conventional
additive sensor noise.

Furthermore, the side-looking acquisition geometry can produce terrain
distortions such as \textit{foreshortening}, \textit{layover}, and
\textit{radar shadow}. SAR products therefore commonly require
radiometric calibration, speckle reduction and terrain correction before
being used in subsequent analyses.

An important example of a spaceborne SAR mission is Copernicus
Sentinel-1, whose instrument operates in the C-band. SAR data are widely
used for flood mapping, soil-moisture analysis, forest monitoring,
sea-ice observation, maritime surveillance and the measurement of
surface deformation. In particular, Interferometric SAR (InSAR) compares
the phase of two or more SAR acquisitions to detect changes in the
distance between the sensor and the Earth's surface.

\subsubsection{LiDAR}

Light Detection and Ranging (LiDAR) is an active optical remote-sensing
technology used to measure the three-dimensional structure of the
Earth's surface. A LiDAR instrument emits short laser pulses toward a
target and records the reflected energy returning to the sensor
\cite{baltsavias1999airborne}. The acquisition process can be represented
as

\[
\text{LiDAR sensor}
\longrightarrow
\text{Earth's surface}
\longrightarrow
\text{LiDAR sensor}.
\]

The distance \(R\) between the sensor and the target is calculated from
the round-trip travel time of the laser pulse:

\[
R = \frac{c\Delta t}{2},
\]

where \(c\) is the speed of light and \(\Delta t\) is the elapsed time
between the transmission and reception of the pulse. The measured range
is combined with the position and orientation of the sensor, generally
obtained using Global Navigation Satellite Systems (GNSS) and an
Inertial Navigation System (INS), to determine the three-dimensional
coordinates of the reflecting point.

By repeatedly emitting laser pulses while the sensor moves, LiDAR
systems produce a collection of three-dimensional points known as a
\textit{point cloud}. A LiDAR point may be represented as

\[
\mathbf{p}_i =
\left[
x_i,\,
y_i,\,
z_i,\,
I_i,\,
r_i
\right],
\]

where \(x_i\), \(y_i\), and \(z_i\) are its spatial coordinates,
\(I_i\) is the intensity of the returned signal, and \(r_i\) identifies
its return number.

A single laser pulse can produce multiple returns when it encounters
objects at different heights. In a forest, for example, the first
returns may originate from the upper canopy, while later returns may
originate from branches, understory vegetation, or the ground. Ground
measurements are obtained when part of the laser energy passes through
gaps in the canopy rather than directly penetrating the vegetation.

LiDAR systems can be deployed on satellites, aircraft, unmanned aerial
vehicles, vehicles, or terrestrial platforms. Their measurements can be
used to generate Digital Terrain Models (DTMs), Digital Surface Models
(DSMs), canopy-height models, and three-dimensional representations of
buildings and infrastructure. LiDAR is therefore widely applied to
topographic mapping, forestry, urban modelling, archaeology, coastal
monitoring, and natural-hazard assessment.

The quality of LiDAR data depends on factors such as point density,
laser-footprint size, scanning geometry, platform position and
orientation, surface reflectivity, and atmospheric conditions. In
addition, conventional topographic LiDAR is affected by clouds and
generally provides limited measurements over water because much of the
laser energy is reflected away from the sensor or absorbed
\cite{usgs_lidar_data}.

\subsubsection{Digital Elevation Models}

A Digital Elevation Model (DEM) is a digital representation of the
elevation of the Earth's surface. It is commonly stored as a
georeferenced raster in which each pixel contains an elevation value:

\[
D(i,j) = z_{i,j},
\]

where \(z_{i,j}\) represents the elevation associated with the spatial
location of pixel \((i,j)\).

A DEM is not itself a remote-sensing sensor. Instead, it is a derived
geospatial product that may be generated from LiDAR point clouds,
stereoscopic optical imagery, Synthetic Aperture Radar interferometry,
photogrammetry, satellite altimetry, or digitised topographic maps
\cite{guth2021dem}.

The term DEM is sometimes used broadly to describe different types of
digital elevation surfaces. More specific terms distinguish the surface
represented by the elevation values:

\begin{itemize}
    \item a \textit{Digital Surface Model} (DSM) represents the elevation
    of the uppermost observed surface, including buildings, vegetation,
    and other objects;

    \item a \textit{Digital Terrain Model} (DTM) represents the underlying
    bare-earth terrain after buildings and vegetation have been removed.
\end{itemize}

In some operational contexts, including that adopted by the United
States Geological Survey, the term DEM is used specifically for a
bare-earth elevation model. The intended meaning should therefore be
stated explicitly whenever an elevation dataset is introduced
\cite{usgs_dem_definition}.

The difference between a DSM and a DTM can also be used to calculate a
normalised Digital Surface Model:

\[
\mathrm{nDSM}(i,j)
=
\mathrm{DSM}(i,j)-\mathrm{DTM}(i,j).
\]

The resulting values represent the approximate height of objects above
the terrain and can therefore be used to estimate building heights,
vegetation heights, and forest-canopy structure.

Several topographic variables can be derived from a DEM, including
slope, aspect, curvature, drainage direction, watershed boundaries, and
terrain visibility. Elevation models are consequently important for
hydrological modelling, flood-risk analysis, landslide assessment,
geomorphology, urban planning, orthorectification, and the geometric
correction of satellite imagery.

When used in machine-learning models, a DEM can be treated as an
additional input channel:

\[
\mathbf{x}_p =
\left[
B_1(p),\ldots,B_n(p),D(p)
\right],
\]

where \(B_1,\ldots,B_n\) are optical or radar measurements and \(D(p)\)
is the elevation associated with pixel \(p\). Terrain information may
help a model distinguish classes or processes that have similar
spectral responses but occur at different elevations, slopes, or
orientations.

The spatial resolution of a DEM describes the horizontal dimensions of
its grid cells and must not be confused with its vertical accuracy.
DEM quality also depends on the original acquisition technology,
interpolation procedure, terrain characteristics, coordinate reference
system, vertical datum, acquisition date, and the presence of voids or
artefacts. These properties must be considered before combining an
elevation model with other Earth-observation data.
%------------------------------------------------

\subsection{Remote Sensing Resolutions}

Remote-sensing systems are commonly characterised through four main
types of resolution:

\begin{itemize}
    \item \textit{Spatial resolution} describes the ground area represented
    by each image pixel. For example, a spatial resolution of
    \(10\,\mathrm{m}\) means that one pixel represents an area of
    approximately \(10 \times 10\,\mathrm{m}\). A smaller pixel size generally
    allows finer spatial details to be distinguished.

    \item \textit{Spectral resolution} describes the ability of a sensor to
    distinguish different wavelength intervals. It depends on the number
    and width of the acquired spectral bands. Sensors with many narrow
    bands generally have higher spectral resolution than sensors with a
    few broad bands.

    \item \textit{Temporal resolution} represents how frequently a sensor
    can observe the same area of the Earth's surface. It is commonly
    expressed through the revisit time. A high temporal resolution is
    particularly important for monitoring rapidly changing phenomena
    such as floods, fires, crops, and cloud cover.

    \item \textit{Radiometric resolution} describes the sensitivity of a
    sensor to differences in the recorded electromagnetic energy. It is
    usually expressed as a number of bits. A sensor with \(n\)-bit
    radiometric resolution can represent

    \[
    2^n
    \]

    different intensity levels. For example, a 12-bit sensor can
    represent \(4096\) possible values.
\end{itemize}

These resolutions are not independent, and improving one of them may
require compromises in the others because of limitations related to
sensor design, data volume, signal strength, and acquisition geometry.


%================================================

\subsection{EO Data Sources and Access}

The practical accessibility of Earth Observation data is an important
consideration for AI4EO research. In addition to the physical
characteristics of a sensor, the feasibility and reproducibility of a
machine-learning study depend on whether the corresponding data can be
searched, accessed, and processed programmatically at the required
spatial and temporal scale. Open archives and standardized interfaces
are therefore particularly valuable when constructing training and
evaluation datasets.

Table~\ref{tab:eo-data-access} summarizes representative data-access
ecosystems that are relevant to AI4EO research.

\begin{table}[htbp]
    \centering
    \small
    \begin{tabular}{
        p{0.18\textwidth}
        p{0.22\textwidth}
        p{0.20\textwidth}
        p{0.16\textwidth}
        p{0.12\textwidth}}
        \toprule
        \textbf{Data Source} &
        \textbf{Representative Data} &
        \textbf{Programmatic Access} &
        \textbf{Interactive Access} &
        \textbf{Access Model} \\
        \midrule

        Copernicus Data Space Ecosystem &
        Sentinel-1, Sentinel-2, Sentinel-3, Sentinel-5P and other Copernicus collections &
        STAC, OData, Sentinel Hub and other APIs &
        Copernicus Browser &
        Open \\

        USGS / EROS &
        Landsat and other EROS archive products &
        Machine-to-Machine (M2M) API &
        EarthExplorer &
        Open \\

        NASA Earthdata &
        MODIS, ICESat-2, SMAP and other NASA Earth-science products &
        CMR-based services and \texttt{earthaccess} &
        Earthdata Search &
        Open \\

        Planet &
        PlanetScope and other Planet imagery products &
        Data API and Processing API &
        Planet data applications &
        Commercial / account-dependent \\

        \bottomrule
    \end{tabular}
    \caption{Representative Earth Observation data-access ecosystems relevant to AI4EO research.}
    \label{tab:eo-data-access}
\end{table}

\paragraph{Data-Access Implementations}

Python clients previously developed for programmatic access
to CDSE and the USGS M2M API are publicly available at:

\begin{itemize}
    \item \url{https://github.com/stefanoIn/CDSE-api}
    \item \url{https://github.com/stefanoIn/USGS-m2m-client}
\end{itemize}

\subsubsection{Copernicus Data Space Ecosystem}

The Copernicus Data Space Ecosystem (CDSE) provides access to data from
the Copernicus Sentinel missions and additional Earth Observation
collections \cite{cdse_documentation}. The Copernicus Browser supports
interactive discovery, visualization, and download of EO products,
whereas programmatic workflows can use interfaces such as the STAC
Catalogue API, OData, and Sentinel Hub services
\cite{cdse_apis,cdse_browser}.

These interfaces are particularly useful in AI4EO because training and
evaluation datasets can be constructed systematically using spatial,
temporal, sensor, product-level, and cloud-cover constraints rather than
through manual scene selection. They also support workflows in which
large EO archives are queried or processed directly from software
pipelines.

\subsubsection{USGS EarthExplorer and Machine-to-Machine API}

The United States Geological Survey distributes Landsat and other Earth
Observation products through the EROS archive. EarthExplorer provides an
interactive interface for discovering and accessing these products,
while the Machine-to-Machine (M2M) API provides a programmatic method for
searching and acquiring data from the same inventories
\cite{usgs_m2m}.

For AI4EO experiments, programmatic access makes it possible to define
reproducible acquisition procedures based on geographic regions,
acquisition dates, products, and metadata filters. This is particularly
useful when constructing large multi-temporal Landsat datasets or when
the same data-selection procedure must be repeated across different
study areas.

\subsubsection{NASA Earthdata}

NASA Earthdata provides access to a broad range of Earth-science
products, including satellite, airborne, and model-derived datasets.
Earthdata Search supports interactive discovery and access, while NASA
also provides software and APIs for programmatic workflows
\cite{nasa_earthdata_search}.

The \texttt{earthaccess} Python library simplifies authentication,
search, download, and cloud-based access to NASA Earth-science data
\cite{nasa_earthaccess}. Such tools are relevant when AI4EO experiments
require repeated access to large collections or integration with
cloud-based processing environments.

\subsubsection{Commercial EO Data}

Commercial providers can offer very-high-resolution imagery, high
revisit frequencies, or specialized products that are not available
through fully open public archives. Planet, for example, provides
programmatic catalogue search and imagery access through its Data API
and Processing API \cite{planet_data_api,planet_processing_api}.

Commercial data can be scientifically valuable, but licensing,
coverage, quotas, and acquisition costs may restrict reproducibility or
make a research direction dependent on institutional access. These
constraints should therefore be considered explicitly when evaluating
the feasibility of a thesis project.

%================================================
\section{AI Tasks in Earth Observation}

This section focuses on the main problem formulations encountered in
AI4EO. The task categories describe the type of mapping learned from EO
data rather than a particular application domain. A single application,
such as wildfire assessment, agriculture, or environmental monitoring,
may therefore combine several task families. The emphasis is placed on
EO-specific inputs, outputs, and application requirements rather than on
reintroducing standard machine-learning or computer-vision algorithms.

\subsection{Image Classification}

Image classification assigns a semantic label to an entire Earth
Observation image or image patch. Given an input image

\[
\mathbf{X}
\in
\mathbb{R}^{C\times H\times W},
\]

the task can be formulated as

\[
f:
\mathbb{R}^{C\times H\times W}
\rightarrow
\mathbb{N}.
\]

For a classification problem with \(K\) possible classes, the output can
be written more precisely as

\[
y \in \{1,\ldots,K\}.
\]

In Earth Observation, image classification is commonly used for scene
classification and for assigning a single dominant semantic category to
an image patch, parcel, or region. Crop-type or land-use labels can be
formulated in this way when one label is assigned to the whole spatial
unit; dense pixel-level land-cover mapping instead falls under semantic
segmentation. The AID benchmark, for example, contains aerial images
assigned to scene categories such as farmland, forest, industrial areas,
residential areas, and airports \cite{xia2017aid}.

\begin{figure}[!htbp]
    \centering
    \includegraphics[
        width=\linewidth,
        keepaspectratio
    ]{images/aid_classification.png}
    \caption{
        Example images from the AID dataset, illustrating the aerial
        scene-classification task. Reproduced from \cite{peng2022mopnet}.
    }
    \label{fig:aid-classification}
\end{figure}

\subsection{Image Captioning}

Image captioning generates a natural-language description of an Earth
Observation image. Given an input image

\[
\mathbf{X}
\in
\mathbb{R}^{C\times H\times W},
\]

the task produces a sequence of \(L\) tokens describing the content of
the observed scene. For a vocabulary containing \(V\) possible tokens,
the task can be formulated as

\[
f:
\mathbb{R}^{C\times H\times W}
\rightarrow
\{1,\ldots,V\}^{L}.
\]

In Earth Observation, image captioning can be used to describe land-cover
patterns, objects, spatial relationships, and other relevant properties
of satellite or aerial imagery.

\begin{figure}[!htbp]
    \centering
    \includegraphics[
        width=\linewidth,
        keepaspectratio
    ]{images/image_captioning.jpg}
    \caption{
        Examples of captions generated for remote-sensing images by
        different image-captioning frameworks. Reproduced from
        \cite{ni2024objectcounts}.
    }
    \label{fig:eo-image-captioning}
\end{figure}


\subsection{Semantic Segmentation}

Semantic segmentation assigns a semantic class to every pixel of an
Earth Observation image. Given an input image

\[
\mathbf{X}
\in
\mathbb{R}^{C\times H\times W},
\]

the task can be formulated as

\[
f:
\mathbb{R}^{C\times H\times W}
\rightarrow
\mathbb{N}^{H\times W},
\]

where each output value identifies the class associated with the
corresponding spatial location.

For a problem with \(K\) semantic classes, the output can be written more
precisely as

\[
\mathbf{Y}
\in
\{1,\ldots,K\}^{H\times W}.
\]

Semantic segmentation is widely used in Earth Observation for tasks such
as land-cover mapping, flood delineation, burned-area mapping, building
segmentation, road extraction, and vegetation mapping. Compared with
standard image classification, it preserves the spatial structure of the
scene by producing a dense pixel-level prediction.

\begin{figure}[!htbp]
    \centering
    \includegraphics[
        width=\linewidth,
        keepaspectratio
    ]{images/semantic_segmentation.png}
    \caption{
        Example of semantic segmentation applied to aerial imagery,
        showing the input image and the corresponding pixelwise semantic
        mask. Adapted from \cite{davies2022semantic}.
    }
    \label{fig:semantic-segmentation}
\end{figure}

\subsection{Instance Segmentation}

Instance segmentation identifies and separates individual object
instances within an Earth Observation image. Given an input image

\[
\mathbf{X}
\in
\mathbb{R}^{C\times H\times W},
\]

the model produces a set of \(D\) spatial masks, which can be represented
as

\[
f:
\mathbb{R}^{C\times H\times W}
\rightarrow
\{0,1\}^{D\times H\times W},
\]

where \(D\) denotes the number of detected object instances, which is
generally image-dependent, and each \(H\times W\) mask identifies the
pixels belonging to one instance.

Unlike semantic segmentation, which assigns pixels only to semantic
classes, instance segmentation distinguishes different objects belonging
to the same class. In Earth Observation, it can be used for applications
such as individual building, vehicle, ship, or tree-crown delineation.

\begin{figure}[!htbp]
    \centering
    \includegraphics[
        width=\linewidth,
        keepaspectratio
    ]{images/instance_segmentation.png}
    \caption{
        Examples illustrating challenges of instance segmentation in
        remote-sensing imagery, including scale variation, low contrast,
        arbitrary orientation, clustered object distributions, and
        cluttered backgrounds. Reproduced from \cite{liu2025catnet}.
    }
    \label{fig:instance-segmentation}
\end{figure}


\subsection{Object Detection}

Object detection aims to identify and localize individual objects within
an Earth Observation image. Given an input image

\[
\mathbf{X}
\in
\mathbb{R}^{C\times H\times W},
\]

the localization component of the task can be represented as

\[
f:
\mathbb{R}^{C\times H\times W}
\rightarrow
\mathbb{R}^{D\times 4},
\]

where \(D\) denotes the number of detected objects and each row contains
the coordinates defining one bounding box.

In practice, each detected object is also associated with a semantic
class and usually a confidence score. Object detection is particularly
relevant in very-high-resolution EO imagery for detecting objects such
as buildings, vehicles, ships, aircraft, and infrastructure.

\begin{figure}[!htbp]
    \centering
    \includegraphics[
        width=\linewidth,
        keepaspectratio
    ]{images/object_detection.jpg}
    \caption{
        Example of object detection in high-resolution remote-sensing
        imagery, where individual objects are localized using bounding
        boxes. Reproduced from \cite{ballinger2023rs4osint}.
    }
    \label{fig:object-detection}
\end{figure}



\subsection{Change Detection}

Change detection aims to identify areas of the Earth's surface that have
changed between two or more observations acquired at different times.
In the bi-temporal case, two co-registered Earth Observation images
acquired at times \(t_1\) and \(t_2\) can be represented as

\[
\mathbf{X}_{t_1},\mathbf{X}_{t_2}
\in
\mathbb{R}^{C\times H\times W},
\]

where \(C\) is the number of input channels and \(H\) and \(W\) are the
spatial dimensions.

The two observations can be considered jointly as an input tensor

\[
\mathbf{X}
\in
\mathbb{R}^{2\times C\times H\times W}.
\]

For binary change detection, the task can be formulated as

\[
f:
\mathbb{R}^{2\times C\times H\times W}
\rightarrow
\{0,1\}^{H\times W},
\]

where each output pixel indicates whether a change has occurred at the
corresponding spatial location:

\[
Y_{i,j}=
\begin{cases}
0, & \text{no change},\\
1, & \text{change}.
\end{cases}
\]

More generally, semantic change detection can distinguish among multiple
types of change and may therefore be written as

\[
f:
\mathbb{R}^{2\times C\times H\times W}
\rightarrow
\{1,\ldots,K\}^{H\times W},
\]

where \(K\) represents the number of possible change classes.

A major challenge in Earth Observation change detection is distinguishing
actual changes on the ground from differences caused by acquisition
conditions. Variations in illumination, atmospheric conditions,
seasonality, viewing geometry, sensor characteristics, spatial
misregistration, and cloud cover may produce apparent differences
between acquisitions even when no meaningful surface change has occurred.

Change detection is widely used for applications such as urban
expansion, deforestation, flood mapping, wildfire assessment,
infrastructure monitoring, agricultural monitoring, and disaster
response. Depending on the application, the observations may originate
from the same sensor or from different sensing modalities, such as
optical and SAR imagery.


\subsection{Regression and Geophysical Parameter Estimation}
Regression tasks estimate continuous-valued quantities from Earth
Observation data. Depending on the application, the prediction may be
associated with an entire image or spatial region,

\[
f:\mathbb{R}^{C\times H\times W}\rightarrow\mathbb{R},
\]

or may consist of a spatially distributed continuous field,

\[
f:\mathbb{R}^{C\times H\times W}
\rightarrow
\mathbb{R}^{H\times W}.
\]

The latter formulation is commonly referred to as pixelwise regression.
Typical applications include biomass estimation, land-surface
temperature estimation, soil-moisture retrieval, and other biophysical
or geophysical parameter-estimation problems.

\subsection{Time-Series Analysis and Forecasting}

Satellite image time series consist of repeated observations of the same
geographical area acquired over time. Given a sequence of \(T\) Earth
Observation images,

\[
\mathbf{X}
\in
\mathbb{R}^{T\times C\times H\times W},
\]

the temporal dimension can be exploited for different downstream
objectives. Depending on the application, the output may correspond to a
class, a continuous variable, a spatial map, a detected event or change,
or a future state. Time-series analysis therefore overlaps with several
of the task formulations introduced above, but differs in that the
sequence of observations is explicitly modelled rather than treating
acquisitions independently.

For example, spatial forecasting can be formulated as

\[
f:
\mathbb{R}^{T\times C\times H\times W}
\rightarrow
\mathbb{R}^{C'\times H\times W},
\]

where the output represents one or more spatially distributed variables
at a future time.

Time-series methods are relevant to applications such as vegetation and
crop monitoring, air-quality monitoring, land-cover and environmental
change analysis, disaster monitoring, and weather and climate
forecasting. Important EO-specific challenges include irregular
acquisition times, missing observations, cloud contamination, seasonal
variability, long temporal dependencies, and differences in temporal
resolution between sensors.

\subsection{Image Reconstruction and Enhancement}

Image reconstruction and enhancement tasks aim to recover or improve an
Earth Observation image rather than directly predicting semantic labels.
In general, they can be represented as

\[
f:
\mathbb{R}^{C_{\mathrm{in}}\times H_{\mathrm{in}}\times W_{\mathrm{in}}}
\rightarrow
\mathbb{R}^{C_{\mathrm{out}}\times H_{\mathrm{out}}\times W_{\mathrm{out}}},
\]

where the output is another continuous-valued image or EO product.

The exact formulation depends on the application. For example, denoising,
cloud removal, and gap filling typically preserve the spatial dimensions,

\[
f:
\mathbb{R}^{C\times H\times W}
\rightarrow
\mathbb{R}^{C\times H\times W},
\]

whereas super-resolution increases the spatial resolution,

\[
f:
\mathbb{R}^{C\times H\times W}
\rightarrow
\mathbb{R}^{C\times sH\times sW},
\]

with \(s\) denoting the spatial upscaling factor.

Typical EO applications include super-resolution, pansharpening, image
fusion, denoising, cloud removal, and reconstruction of missing
observations.


%================================================
\section{Evolution of AI for Earth Observation}

This section should provide only the EO-specific historical transition
from hand-crafted remote-sensing pipelines to learned representations.
Standard AI architectures need not be explained in tutorial form.

\subsection{From Hand-Crafted Features to Learned Representations}

% Briefly discuss spectral indices, texture/shape features, dimensionality
% reduction, and classical classifiers only as historical context.
% The key transition is from domain-designed features to end-to-end learned
% spectral, spatial, and spectral--spatial representations.

\subsection{Deep Learning for EO}

% Summarize what deep learning changed in EO: joint spectral--spatial learning,
% dense prediction, end-to-end feature learning, and large-scale image processing.
% Mention CNN/FCN/U-Net/ResNet only when they are relevant to the EO transition.

\subsection{Representation Learning and Self-Supervision}

% Explain why unlabeled EO archives make self-supervised learning particularly
% relevant; include contrastive and masked-image/modeling approaches.

\subsection{Transformers and Large-Scale Pretraining}

% Focus on EO-specific motivations: long-range spatial context, multi-temporal
% data, heterogeneous bands/modalities, and scaling to large pretraining corpora.

\subsection{The Shift Toward Earth Observation Foundation Models}

% Introduce the move from task-specific models to reusable pretrained models.
% Detailed comparison of current foundation models belongs in the next section.

%================================================
\section{Current Research Trends}

This section forms the core of the report. The objective is to identify
recurring research problems and methodological directions across the
recent AI4EO literature rather than to infer the importance of a topic
from a single workshop or conference. In particular, workshops focused
explicitly on Earth Observation foundation models are useful for
understanding developments within that research branch, but should not
be treated as an unbiased representation of AI4EO as a whole.

Each trend should therefore be supported by representative recent papers
from multiple venues and should summarize the underlying research
problem, current approaches, main results, and remaining limitations.

\subsection{Large-Scale Pretraining and Earth Observation Foundation Models}

% Discuss the transition from task-specific EO models to reusable pretrained
% representations. Representative themes include self-supervised pretraining,
% masked modelling, sensor-specific versus sensor-agnostic models,
% multimodal/multiresolution foundation models, downstream adaptation, and
% model evaluation. Representative systems can include Prithvi, TerraMind,
% AnySat, PhilEO, DOFA, SatMAE, CROMA, and other models supported by the
% literature review. Avoid treating FM-focused workshops as evidence that
% foundation models constitute the entirety of current AI4EO research.

\subsection{Multimodal and Multisensor Learning}

% Optical + SAR + hyperspectral + LiDAR/DEM + weather/climate variables +
% in-situ observations + text or other geospatial information.
% Discuss feature-level, representation-level and decision-level fusion,
% modality-specific preprocessing, co-registration, asynchronous acquisitions,
% resolution mismatch, and missing modalities.

\subsection{Spatio-Temporal Modelling and Change Analysis}

% Satellite image time series, bi-temporal and multi-temporal change detection,
% forecasting, event evolution, irregular sampling, cloud-related gaps,
% seasonality, long temporal context, and distinguishing semantic change from
% acquisition or phenological variation.

\subsection{Generalization, Domain Adaptation, and Limited Supervision}

% Geographic, seasonal, sensor and resolution shifts; transfer learning,
% domain adaptation, semi-supervised learning, weak supervision, few-shot
% learning, incomplete annotations, active learning, and out-of-distribution
% detection. Emphasize the common EO problem: labelled training data often do
% not match the deployment geography, season, sensor, or acquisition setting.

\subsection{Trustworthy, Explainable, and Uncertainty-Aware AI}

% Uncertainty quantification, calibration, explainability, robustness,
% out-of-distribution behaviour, traceability, failure detection, and
% scientific reliability. Emphasize whether predictions remain trustworthy
% across geography, season, sensor, and acquisition conditions.

\subsection{Physics-Aware and Hybrid AI}

% Physically meaningful constraints, differentiable physical models,
% hybrid physical/data-driven approaches, conservation laws, radiative-transfer
% knowledge, and scientific priors. Distinguish this branch from generic
% explainability and reliability research.

\subsection{Vision-Language and Agentic Geospatial AI}

% Image--text representation learning, vision-language foundation models,
% retrieval, captioning, grounding, natural-language interaction with EO data,
% domain-specific LLMs, RAG, tool use, and agentic workflows. Treat this as an
% emerging research direction and distinguish visual-language representation
% learning from LLM-based scientific assistants.

\subsection{Efficient, Scalable, and Operational AI}

% Parameter-efficient adaptation, distillation, compression, efficient
% inference, cloud/HPC deployment, large-scale serving, learned data
% compression, edge processing, onboard AI, and constraints imposed by
% satellite hardware and operational workflows.

\subsection{Datasets, Benchmarks, Embeddings, and Representation Infrastructure}

% Dataset construction, standardized benchmarks, geographic train/test design,
% capability-oriented evaluation, pretrained EO embeddings, global embedding
% products, analysis-ready data, and software infrastructure. Representative
% initiatives identified in the literature may include PANGAEA, GEO-Bench,
% GEO-Bench-2, GFM-Bench, EarthNets, FORTY, MDAS, and Earth2Vec-style
% embedding initiatives.

%================================================
\section{Open Challenges and Feasibility}

This section translates the research landscape into the practical and
scientific considerations needed to evaluate possible MSc thesis
directions. The objective is not to select the thesis in advance, but
to identify where important open problems intersect with accessible
data, realistic computational requirements, and clear evaluation
methodologies.

\subsection{Data Availability and Annotation}

% Open versus commercial data, labelled benchmarks, spatial and temporal
% coverage, annotation cost, weak or noisy labels, data imbalance, and
% reproducibility. Connect this subsection with the data-access ecosystems
% described in Section 2.

\subsection{Distribution Shift and Generalization}

% Geographic transfer, seasonal variation, sensor and resolution changes,
% acquisition-condition shifts, robustness outside the training distribution,
% and the distinction between benchmark performance and deployment reliability.

\subsection{Multimodal and Multi-Resolution Integration}

% Co-registration, asynchronous acquisition, missing modalities, heterogeneous
% native resolutions, resampling, differing physical measurement processes,
% and modality-specific preprocessing.

\subsection{Computational and Infrastructure Requirements}

% Pretraining versus fine-tuning, GPU memory, storage, data-transfer costs,
% data pipelines, training time, model size, cloud/HPC availability, and
% whether existing pretrained models or embeddings make a direction feasible.

\subsection{Evaluation, Benchmarking, and Reproducibility}

% Benchmark quality, geographic train/test separation, temporal leakage,
% sensor leakage, appropriate metrics, baselines, statistical reliability,
% reproducibility, and whether a proposed contribution can be evaluated with a
% sufficiently clear experimental protocol.

\subsection{Candidate Research Directions}

Candidate directions identified from the literature can be compared using
criteria such as the underlying research gap, scientific relevance,
availability of data, computational requirements, implementation complexity,
and evaluation methodology. The final thesis direction can then be selected
jointly based on these trade-offs.

\begin{table}[htbp]
    \centering
    \small
    \begin{tabular}{
        p{0.18\textwidth}
        p{0.11\textwidth}
        p{0.11\textwidth}
        p{0.14\textwidth}
        p{0.17\textwidth}
        p{0.17\textwidth}}
        \toprule
        \textbf{Direction} &
        \textbf{Data} &
        \textbf{Compute} &
        \textbf{Evaluation} &
        \textbf{Research Gap} &
        \textbf{Feasibility} \\
        \midrule
        -- & -- & -- & -- & -- & -- \\
        \bottomrule
    \end{tabular}
    \caption{Template for comparing candidate research directions.}
    \label{tab:candidate-directions}
\end{table}

%================================================
\section{Conclusion}

% Concisely summarize the state of AI4EO, the strongest research trends,
% and the principal feasibility constraints. Do not preselect the thesis topic.


\begin{thebibliography}{99}

\bibitem{esa_optical_imagery}
European Space Agency,
\textit{Optical Earth Observation},
ESA Earth Observation.
Available online:
\url{https://www.esa.int/Applications/Observing_the_Earth}

\bibitem{esa_reflectance_curves}
European Space Agency,
\textit{Reflectance Curves of Snow, Vegetation, Water and Rock},
ESA Eduspace.
Available online:
\url{https://www.esa.int/SPECIALS/Eduspace_EN/}

\bibitem{esa_sentinel2_msi}
European Space Agency,
\textit{Sentinel-2 Multispectral Instrument Overview},
ESA Sentinel Application Platform.
Available online:
\url{https://step.esa.int/main/wp-content/help/versions/11.0.0/snap-toolboxes/eu.esa.opt.opttbx.s2msi.reader/Sentinel2Overview.html}

\bibitem{jpl_imaging_spectroscopy}
NASA Jet Propulsion Laboratory,
\textit{AVIRIS Imaging Spectroscopy}.
Available online:
\url{https://aviris.jpl.nasa.gov/aviris/imaging_spectroscopy.html}

\bibitem{bioucasdias2013hyperspectral}
J.~M. Bioucas-Dias, A. Plaza, G. Camps-Valls, P. Scheunders,
N.~M. Nasrabadi, and J. Chanussot,
``Hyperspectral remote sensing data analysis and future challenges,''
\textit{IEEE Geoscience and Remote Sensing Magazine},
vol.~1, no.~2, pp.~6--36, 2013.

\bibitem{rouse1974monitoring}
J.~W. Rouse, R.~H. Haas, J.~A. Schell, and D.~W. Deering,
``Monitoring vegetation systems in the Great Plains with ERTS,''
in \textit{Proceedings of the Third Earth Resources Technology
Satellite-1 Symposium},
vol.~1, pp.~309--317, 1974.

\bibitem{moreira2013sar}
A.~Moreira, P.~Prats-Iraola, M.~Younis, G.~Krieger, I.~Hajnsek,
and K.~P. Papathanassiou,
``A tutorial on synthetic aperture radar,''
\textit{IEEE Geoscience and Remote Sensing Magazine},
vol.~1, no.~1, pp.~6--43, 2013.

\bibitem{esa_sentinel1_instrument}
European Space Agency,
\textit{Sentinel-1 Instrument},
ESA Earth Observation.
Available online:
\url{https://www.esa.int/Applications/Observing_the_Earth/Copernicus/Sentinel-1/Instrument}

\bibitem{baltsavias1999airborne}
E.~P. Baltsavias,
``Airborne laser scanning: Basic relations and formulas,''
\textit{ISPRS Journal of Photogrammetry and Remote Sensing},
vol.~54, no.~2--3, pp.~199--214, 1999.

\bibitem{usgs_lidar_data}
United States Geological Survey,
\textit{What Is Lidar Data and Where Can I Download It?}
Available online:
\url{https://www.usgs.gov/faqs/what-lidar-data-and-where-can-i-download-it}

\bibitem{guth2021dem}
P.~L. Guth, A. Van Niekerk, C.~H. Grohmann, J.-P. Muller,
L. Hawker, I.~V. Florinsky, D.~B. Gesch, H.~I. Reuter,
V. Herrera-Cruz, S. Riazanoff, C. Lopez-Vazquez,
C.~C. Carabajal, C. Albinet, and P. Strobl,
``Digital elevation models: Terminology and definitions,''
\textit{Remote Sensing},
vol.~13, no.~18, Art.~no.~3581, 2021.

\bibitem{usgs_dem_definition}
United States Geological Survey,
\textit{What Is a Digital Elevation Model (DEM)?}
Available online:
\url{https://www.usgs.gov/faqs/what-a-digital-elevation-model-dem}

\bibitem{esa_earth_observation}
European Space Agency,
\textit{Newcomers Earth Observation Guide}.
Available online:
\url{https://business.esa.int/newcomers-earth-observation-guide}


\bibitem{cdse_documentation}
Copernicus Data Space Ecosystem,
\textit{Copernicus Data Space Ecosystem Documentation}.
Available online:
\url{https://documentation.dataspace.copernicus.eu/}

\bibitem{cdse_apis}
Copernicus Data Space Ecosystem,
\textit{APIs}.
Available online:
\url{https://documentation.dataspace.copernicus.eu/APIs.html}

\bibitem{cdse_browser}
Copernicus Data Space Ecosystem,
\textit{About the Copernicus Browser}.
Available online:
\url{https://documentation.dataspace.copernicus.eu/Applications/Browser.html}

\bibitem{usgs_m2m}
United States Geological Survey,
\textit{Machine-to-Machine (M2M) API}.
Available online:
\url{https://m2m.cr.usgs.gov/}

\bibitem{nasa_earthdata_search}
National Aeronautics and Space Administration,
\textit{Earthdata Search}.
Available online:
\url{https://www.earthdata.nasa.gov/data/tools/earthdata-search}

\bibitem{nasa_earthaccess}
National Aeronautics and Space Administration,
\textit{earthaccess}.
Available online:
\url{https://www.earthdata.nasa.gov/data/tools/earthaccess}

\bibitem{planet_data_api}
Planet Labs,
\textit{Data API Overview}.
Available online:
\url{https://docs.planet.com/develop/apis/data/}

\bibitem{planet_processing_api}
Planet Labs,
\textit{Processing API}.
Available online:
\url{https://docs.planet.com/develop/apis/processing/}

\bibitem{peng2022mopnet}
F.~Peng, W.~Lu, W.~Tan, K.~Qi, X.~Zhang, and Q.~Zhu,
``Multi-Output Network Combining GNN and CNN for Remote Sensing Scene Classification,''
\textit{Remote Sensing},
vol.~14, no.~6, Art.~no.~1478, 2022.

\bibitem{xia2017aid}
G.-S. Xia, J.~Hu, F.~Hu, B.~Shi, X.~Bai, Y.~Zhong, L.~Zhang,
and X.~Lu,
``AID: A Benchmark Data Set for Performance Evaluation of Aerial
Scene Classification,''
\textit{IEEE Transactions on Geoscience and Remote Sensing},
vol.~55, no.~7, pp.~3965--3981, 2017.

\bibitem{ni2024objectcounts}
Z.~Ni, Z.~Zong, and P.~Ren,
``Incorporating object counts into remote sensing image captioning,''
\textit{International Journal of Digital Earth},
vol.~17, no.~1, Art.~no.~2392847, 2024.

\bibitem{davies2022semantic}
A.~J. Davies,
\textit{Semantic Segmentation of Aerial Imagery using U-Net in Python},
Towards Data Science, Mar.~31, 2022.
Available online:
\url{https://towardsdatascience.com/semantic-segmentation-of-aerial-imagery-using-u-net-in-python-552705238514/}

\bibitem{liu2025catnet}
Y.~Liu, H.~Li, C.~Hu, S.~Luo, Y.~Luo, and C.~W. Chen,
``Learning to Aggregate Multi-Scale Context for Instance Segmentation
in Remote Sensing Images,''
\textit{IEEE Transactions on Neural Networks and Learning Systems},
vol.~36, no.~1, pp.~595--609, 2025.

\bibitem{ballinger2023rs4osint}
O.~Ballinger,
\textit{Remote Sensing for OSINT: Object Detection},
2023.
Available online:
\url{https://bellingcat.github.io/RS4OSINT/C5_Object_Detection.html}

\end{thebibliography}

\end{document}